\section{Introduction}\label{sec:introduction}

Can fiscal space generated by artificial intelligence finance recurrent social protection in Latin American economies---and, if so, when? AI may raise output, alter factor incomes, and expand taxable bases. Social pensions, minimum-income guarantees, and universal transfers, however, create commitments that must be financed repeatedly rather than at a favorable point in time. The relevant policy question is therefore not whether an AI dividend is conceivable. It is whether, under explicit scenarios and for long enough, that dividend can be translated into domestic income, captured by the public sector, and made available after policy costs and debt requirements are recognized. That timing distinction matters for budget design: a large eventual dividend cannot support a current entitlement unless its fiscal transmission and duration are credible. Answering the question comparably across countries and instruments requires a common financing test rather than a forecast attached to a single program.

Each link in that chain is uncertain. Macroeconomic gains depend on which tasks AI affects and how those gains diffuse, while automation can change the distribution of income across tax bases \citep{acemoglu2024simple,acemoglu2019automation}. Fiscal policy can broaden or narrow the distribution of AI gains, but effective capture depends on the existing tax architecture, behavioral responses, and administrative capacity \citep{brollo2024fiscal,korinek2026public,gaspar2016tax}. Exposure is therefore not adoption; additional output is not automatically taxable income; and additional revenue is not necessarily debt-consistent fiscal space. These distinctions motivate a threshold exercise that keeps the technological, domestic, fiscal, and policy-cost objects separate.

Two strands of work come closest to the question. AI--universal-basic-income (UBI) models study particular financing arrangements: \citet{nayebi2025threshold} derives an AI-capability threshold for rent-funded UBI in the United States, while \citet{ito2026ubi} simulates UBI feasibility for Japan with dynamic social-security offsets. They move beyond pure incidence, but retain one country and one instrument. In the policy-modeling tradition, \citet{caamal2022ubi} evaluate UBI in Mexico and \citet{lara2026basic} study a basic-income--flat-tax reform in Spain through country-specific tax-benefit simulations. Those studies identify poverty, inequality, efficiency, or budget effects at a specified policy configuration. None provides the comparative dynamic-fiscal test posed here: when, and in which country--instrument--scenario cells, can a projected AI dividend sustain a recurrent commitment after domestic translation and fiscal capture, on a common denominator, while satisfying debt consistency and a required duration? This paper contributes that test. It complements static incidence analysis rather than replacing it: the object is not how much a transfer redistributes at one point in time, but under which conditions its projected financing source can persist.

The framework implements this comparison as a modular chain. An external AI scenario is first translated into domestic income. The income increment is allocated across labor and nonlabor bases and converted into effective fiscal space through marginal fiscal conversion (MFC). That flow is compared with the harmonized cost of each instrument through the feasibility ratio $V$, with the accounting crossing defined by $V\geq1$. A separate guardrail asks whether the same flow can finance the instrument while meeting the primary-balance correction implied by the maintained debt path. Diagnostic inversions then recover the gain, fiscal conversion, and domestic translation required to reach the boundary. A final inversion, $H^*$, reports the required persistence horizon; it is a derived post-baseline extension, not a forecast date, and changes no calibrated input or baseline classification.

The application covers four economies selected for contrast rather than regional representativeness: Peru, Chile, Colombia, and Mexico. They differ in policy costs, fiscal architectures, and debt space. The instrument set spans categorical, targeted, and broad universal commitments, all priced on the same endpoint-GDP denominator. A common scenario grid changes the technological input, while the fiscal regimes change effective capture without relabeling that input. This separation gives a crossing the same accounting meaning across country--instrument cells and shows whether debt consistency changes the classification. The purpose is not to rank the welfare value of the programs, but to test the sufficiency of a specified financing source under transparent and comparable conditions.

The first result is a deterministic non-crossing under the primary channel-based specification across the non-stress scenarios (low, mid, and frontier). The largest accounting ratio is only $V=0.177$. The gap is joint: domestic translation at the maintained fiscal architecture is insufficient, but reform-level capture at the maintained technological gain is also insufficient. At the maintained prudential buffer, every mid-scenario translation inversion requires $T^{req,H}>1$, and the lowest AI-rent capture requirement is 13.50 times the strongest specified upper support. The evidence therefore supports neither ``translation, not taxation'' nor its inverse. Translation and capture are complementary links, and neither closes the non-stress gap alone. In the primary Monte Carlo, the largest non-stress crossing frequency is 2.36 percent. This is a model-implied frequency across calibrated draws, conditional on stated supports and distributional assumptions, not a predictive probability of real-world financing.

Deterministic crossings appear only under the disruptive scenario, which is a diagnostic stress test rather than a central projection. Chile's guaranteed minimum income (GMI) is the sole conditionally feasible country--instrument result: it crosses in both the ideal and targeting-cost-loaded variants and passes the debt guardrail. Peru's social pension crosses the accounting boundary under stress but fails debt consistency, making it the fragile case rather than a fiscal green light. Colombia and Mexico do not cross in any evaluated cell. These negative findings are informative because every cell faces the same financing condition; they delimit what the stated scenarios can support instead of leaving non-feasibility unmeasured. The persistence extension reinforces the conditional reading by asking how long a shock must endure, not when adoption will occur. The ICT placebo, reduced-form exposure-rescaling diagnostic, channel ablations, and source-substitution exercise further test whether the crossing pattern is specific to the maintained channels.

The policy contribution is consequently narrow but operational. Prospective AI revenues should not be booked as short-run financing for universal social protection, and fiscal reform cannot be treated as a substitute for realizing the domestic productive gain. The threshold framework instead identifies the evidence that would warrant reassessment: realized translation, effective capture, debt consistency, and sufficient persistence. After the next section positions the paper in the literature, Sections~\ref{sec:framework} and~\ref{sec:design} develop and operationalize the test. Section~\ref{sec:results} presents the evidence, and Sections~\ref{sec:discussion} and~\ref{sec:conclusion} draw out its policy implications and conclude.
